% ============================================================================
% LANGENTWURF LE-019: Kalender & Kontakte
% Profil: S (Senioren 65+) | Tier: 2 (Standard)
% ============================================================================

\documentclass[11pt,a4paper]{article}
\usepackage[utf8]{inputenc}
\usepackage[T1]{fontenc}
\usepackage[ngerman]{babel}
\usepackage{geometry}
\usepackage{fancyhdr}
\usepackage{lastpage}
\usepackage{xcolor}
\usepackage{tcolorbox}
\usepackage{enumitem}
\usepackage{longtable}
\usepackage{hyperref}
\usepackage{amssymb}

\geometry{left=2.5cm, right=2.5cm, top=2.5cm, bottom=2.5cm}
\definecolor{fach}{HTML}{b35c00}
\definecolor{gray}{HTML}{666666}
\definecolor{successgreen}{HTML}{2a7a4b}

\tcbuselibrary{skins,breakable}
\newtcolorbox{scorebox}{ colback=white, colframe=fach, fonttitle=\bfseries, title=Didaktik-Score, breakable }
\newtcolorbox{tippbox}[1][]{ colback=successgreen!5, colframe=successgreen, fonttitle=\bfseries, title=#1, breakable }

\pagestyle{fancy}
\fancyhf{}
\fancyhead[L]{\small\textcolor{gray}{Digitale Grundlagen -- 019 Kalender \& Kontakte}}
\fancyhead[R]{\small\textcolor{gray}{LE Tier 2 / Profil S}}
\fancyfoot[C]{\small\thepage{} / \pageref{LastPage}}
\renewcommand{\headrulewidth}{0.4pt}

\begin{document}

\begin{titlepage}
    \centering
    \vspace*{2cm}
    {\Large\textcolor{gray}{Digitale Grundlagen -- Senioren 65+}}\\[2cm]
    {\Huge\textcolor{fach}{\textbf{Langentwurf}}}\\[0.5cm]
    {\large Tier 2 | Profil S}\\[2cm]
    {\LARGE\textbf{019 -- Kalender \& Kontakte}}\\[0.5cm]
    {\large Modul D: Produktivität}\\[2cm]
    \begin{tabular}{rl}
        \textbf{Zeitrahmen:} & 60--75 Minuten \\
        \textbf{Vorkenntnisse:} & Apps nutzen, Tippen \\
    \end{tabular}
    \vfill
    {\normalsize Stand: \today}
\end{titlepage}

\tableofcontents
\newpage

\section{Bedingungsanalyse}

\subsection{Ausgangssituation}
\begin{itemize}
    \item Papierkalender wird bevorzugt
    \item Kontakte oft nur auf SIM-Karte oder gar nicht digital
    \item Erinnerungsfunktion unbekannt
    \item Doppelte Kontakte durch verschiedene Quellen
    \item ``Ich vergesse meine Termine''
\end{itemize}

\subsection{Typische Probleme}
\begin{itemize}
    \item ``Wie trage ich einen Termin ein?''
    \item ``Kann das Handy mich erinnern?''
    \item ``Wie speichere ich eine Telefonnummer?''
    \item ``Ich habe den Kontakt dreimal drin''
\end{itemize}

\section{Sachanalyse}

\subsection{Kernkonzepte}
\begin{enumerate}
    \item \textbf{Kalender-App:} Digitaler Terminkalender mit Erinnerungen
    \item \textbf{Termin:} Ereignis mit Datum, Uhrzeit, optional Ort
    \item \textbf{Erinnerung:} Benachrichtigung vor dem Termin
    \item \textbf{Kontakte-App:} Digitales Adressbuch
    \item \textbf{Favoriten:} Schnellzugriff auf wichtige Kontakte
\end{enumerate}

\subsection{Alltagsanalogien}
\begin{tabular}{|l|p{8cm}|}
\hline
\textbf{Begriff} & \textbf{Analogie} \\
\hline
Kalender-App & Taschenkalender, der dich anstupst \\
Termin & Eintrag im Kalender (Arzt, Geburtstag) \\
Erinnerung & Wecker, der vor dem Termin klingelt \\
Kontakte-App & Adressbuch, das du immer dabei hast \\
Favoriten & Die wichtigsten Nummern vorne im Telefonbuch \\
\hline
\end{tabular}

\section{Didaktische Analyse}

\subsection{Stellung in der Reihe}
Praktische Alltagsorganisation. Ergänzt Notizen (016) um zeitgebundene Inhalte.

\subsection{Didaktische Reduktion}
\textbf{Ausgeklammert:} Mehrere Kalender, geteilte Kalender, wiederkehrende komplexe Termine.

\textbf{Fokus:} Einfache Termine mit Erinnerung, Kontakte anlegen und nutzen.

\section{Lernziele}

\begin{tabular}{|c|p{7cm}|c|c|}
\hline
\textbf{Nr.} & \textbf{Die Teilnehmer können...} & \textbf{Stufe} & \textbf{Niveau} \\
\hline
FZ1 & einen Termin mit Datum und Uhrzeit eintragen & Anwenden & Basis \\
FZ2 & eine Erinnerung für den Termin einstellen & Anwenden & Basis \\
FZ3 & einen neuen Kontakt anlegen & Anwenden & Basis \\
FZ4 & einen Kontakt direkt anrufen oder anschreiben & Anwenden & Basis \\
FZ5 & wiederkehrende Termine einrichten (z.B. Geburtstage) & Anwenden & Vertiefung \\
\hline
\end{tabular}

\section{Methodische Analyse}

\subsection{Termin eintragen}
\begin{tippbox}[Termin schnell eintragen]
\begin{enumerate}
    \item Kalender-App öffnen
    \item + Symbol antippen
    \item Titel eingeben (z.B. ``Arzt Dr. Müller'')
    \item Datum und Uhrzeit einstellen
    \item Erinnerung wählen (z.B. ``1 Tag vorher'')
    \item Fertig antippen
\end{enumerate}
\end{tippbox}

\section{Verlaufsplanung}

\begin{longtable}{|p{1cm}|p{2cm}|p{5cm}|p{3cm}|p{2cm}|}
\hline
\textbf{Zeit} & \textbf{Phase} & \textbf{Inhalt} & \textbf{TN-Aktivität} & \textbf{Material} \\
\hline
0--10 & Einstieg & ``Wie merken Sie sich Termine?'' Vergleich analog/digital & Austauschen & -- \\
\hline
10--25 & Kalender & App öffnen, Ansichten (Tag, Monat), Navigation & Mitmachen & S.1 \\
\hline
25--40 & Termin & Neuen Termin eintragen, Erinnerung setzen & Selbst üben & S.1 \\
\hline
40--55 & Kontakte & Kontakte-App, neuen Kontakt anlegen, Favoriten & Üben & S.2 \\
\hline
55--70 & Verbinden & Aus Kontakt anrufen/schreiben, Termin mit Ort & Ausprobieren & S.2-3 \\
\hline
\end{longtable}

\section{Lösungen}

\textbf{Termin eintragen:}
\begin{itemize}
    \item Kalender öffnen (buntes Symbol)
    \item + oben rechts antippen
    \item Titel eingeben
    \item ``Beginnt'' antippen $\rightarrow$ Datum und Zeit wählen
    \item ``Hinweis'' antippen $\rightarrow$ z.B. ``1 Tag davor''
    \item ``Hinzufügen'' / ``Fertig''
\end{itemize}

\textbf{Kontakt anlegen:}
\begin{itemize}
    \item Kontakte-App öffnen
    \item + oben rechts
    \item Vorname, Nachname eingeben
    \item Telefon hinzufügen
    \item Optional: E-Mail, Adresse, Geburtstag
    \item ``Fertig''
\end{itemize}

\textbf{Kontakt als Favorit:}
\begin{itemize}
    \item Kontakt öffnen
    \item Nach unten scrollen
    \item ``Zu Favoriten hinzufügen''
    \item Favoriten erscheinen in der Telefon-App ganz oben
\end{itemize}

\textbf{Praktische Termine:}
\begin{itemize}
    \item Arzttermine (mit Adresse!)
    \item Geburtstage (jährlich wiederholen)
    \item Medikamenten-Erinnerungen
    \item Müllabfuhr
\end{itemize}

\section{Didaktik-Score}

\begin{scorebox}
\begin{tabular}{|c|l|c|c|p{3.5cm}|}
\hline
\# & Dimension & Gew. & Score & Notiz \\
\hline
1 & Differenzierung & 15\% & 8/10 & Basis/Vertiefung klar \\
2 & Taxonomie & 10\% & 9/10 & Anwenden dominant \\
3 & Alltagsrelevanz & 10\% & 10/10 & Termine vergessen = Problem \\
4 & Aktivierung & 10\% & 10/10 & Eigene Termine/Kontakte \\
5 & Scaffolding & 10\% & 10/10 & Schritt-für-Schritt-Box \\
6 & Feedback & 10\% & 8/10 & Checkliste vorhanden \\
7 & Sprachliche Klarheit & 10\% & 9/10 & Einfache Sprache \\
8 & Motivation & 10\% & 10/10 & Nie mehr Termine vergessen \\
9 & Barrierefreiheit & 10\% & 9/10 & Große Darstellung \\
10 & Praktikabilität & 5\% & 9/10 & 70 Min gut \\
\hline
\multicolumn{3}{|r|}{\textbf{GESAMT:}} & \textbf{9.2/10} & \textcolor{successgreen}{\textbf{FREIGABE}} \\
\hline
\end{tabular}
\end{scorebox}

\end{document}
